% ch2_answers.tex -- answers to the Quick check questions of Module 8, Chapter 2.
\begin{chapteranswers}{2}{Structure of the Atom}
  % ---- 2.1 the electron
  \item The sharp shadow shows that cathode rays travel in straight lines from
        the cathode. It could not settle the debate, because light (a wave) also
        travels in straight lines and casts sharp shadows. The result fits both
        models.
  \item The beam fades or disappears. The electrons now hit many gas particles.
        They lose energy and scatter, so they cannot travel far in a narrow beam.
  \item The beam is not deflected, so
        $v = \dfrac{E}{B} = \dfrac{\qty{6.0e4}{V.m^{-1}}}{\qty{2.4e-3}{T}}
        = \qty{2.5e7}{m.s^{-1}}$.
        Answer: \final{$v = \qty{2.5e7}{m.s^{-1}}$}.
  \item (a) With $r = \qty{2.09}{cm} = \qty{0.0209}{m}$: first square $B$,
        $(\qty{0.100}{T})^2 = \qty{0.0100}{T^2}$. Then multiply by $r$:
        $\num{0.0100}\times\num{0.0209} = \num{2.09e-4}$. So
        $\dfrac{q}{m} = \dfrac{E}{B^2 r}
        = \dfrac{\qty{2.00e4}{V.m^{-1}}}{(\qty{0.100}{T})^2\times\qty{0.0209}{m}}
        = \dfrac{\num{2.00e4}}{\num{2.09e-4}} = \qty{9.57e7}{C.kg^{-1}}$.\par
        (b) $m = q \div \dfrac{q}{m} = \dfrac{\qty{1.602e-19}{C}}{\qty{9.57e7}{C.kg^{-1}}}
        = \qty{1.67e-27}{kg}$. This is the proton mass on the data sheet
        (\qty{1.673e-27}{kg}), so the ion is a hydrogen ion: a proton.
        Answers: \final{$q/m = \qty{9.57e7}{C.kg^{-1}}$}, \final{$m = \qty{1.67e-27}{kg}$}.
  \item (a) The work done on each electron becomes kinetic energy:
        $K = qV = \qty{1.602e-19}{C}\times\qty{180}{V} = \qty{2.884e-17}{J}$.
        Rearrange $K = \tfrac12 mv^2$:
        $v^2 = \dfrac{2K}{m} = \dfrac{2\times\qty{2.884e-17}{J}}{\qty{9.109e-31}{kg}}
        = \qty{6.332e13}{m^2.s^{-2}}$, so $v = \sqrt{\num{6.332e13}} = \qty{7.96e6}{m.s^{-1}}$.
        Answer: \final{$v = \qty{8.0e6}{m.s^{-1}}$}.\par
        (b) For an electron, $\dfrac{q}{m} = \dfrac{\qty{1.602e-19}{C}}{\qty{9.109e-31}{kg}}
        = \qty{1.759e11}{C.kg^{-1}}$. Rearrange $\dfrac{q}{m} = \dfrac{v}{Br}$:
        $r = \dfrac{v}{(q/m)\,B}
        = \dfrac{\qty{7.96e6}{m.s^{-1}}}{\qty{1.759e11}{C.kg^{-1}}\times\qty{0.90e-3}{T}}
        = \dfrac{\num{7.96e6}}{\num{1.583e8}} = \qty{0.0503}{m}$.
        Answer: \final{$r = \qty{5.0}{cm}$}.
  % ---- 2.2 Millikan
  \item Rearrange $q = \dfrac{mgd}{V}$ for $V$, with $d = \qty{8.0}{mm} = \qty{0.0080}{m}$
        and $q = 3e$:
        $V = \dfrac{mgd}{q} = \dfrac{\qty{2.4e-15}{kg}\times\qty{9.8}{m.s^{-2}}\times\qty{0.0080}{m}}{3\times\qty{1.602e-19}{C}}
        = \dfrac{\qty{1.88e-16}{J}}{\qty{4.806e-19}{C}} = \qty{392}{V}$.
        Answer: \final{$V = \qty{3.9e2}{V}$}.
  \item $n = \dfrac{q}{e} = \dfrac{\qty{4.0e-19}{C}}{\qty{1.602e-19}{C}} = 2.5$.
        A drop can only gain or lose whole electrons, so its charge must be a
        whole number times $e$. Half an electron is impossible, so the
        measurement must contain an error.
  \item Before the change, the electric force ($qE$, up) balanced the weight
        ($mg$, down). Swapping the plates reverses the field, so the electric
        force now points down. Both forces point down, so the net force is
        $2mg$ downwards. The drop speeds up downwards (until air resistance
        makes it fall at a steady speed).
  % ---- 2.3 nuclear atom
  \item Sketch: a circle filled with positive charge (draw $+$ signs all
        through it), with small electrons ($-$) dotted through it. Label
        ``positive charge spread through the whole atom'' and ``electrons''.
        There is no nucleus. The model predicts only tiny deflections. The
        spread-out charge gives only a weak force. The light electrons cannot
        turn a heavy alpha.
  \item The atom is mostly empty space, so most alphas pass through without
        coming near a nucleus. The nucleus is tiny, so only a very few alphas
        head almost straight at it. For those few, the positive charge is so
        concentrated that the repulsion is huge, and they are turned back.
  \item Use Coulomb's law with $q_1 = 2e$ and $q_2 = 47e$:
        \begin{align*}
          F &= \frac{1}{4\pi\varepsilon_0}\frac{q_1 q_2}{r^2}
          = \frac{\qty{8.99e9}{N.m^2.C^{-2}}\times(\qty{3.204e-19}{C})\times(\qty{7.529e-18}{C})}{(\qty{2.5e-14}{m})^2}\\
          &= \frac{\qty{8.99e9}{N.m^2.C^{-2}}\times\qty{2.412e-36}{C^2}}{\qty{6.25e-28}{m^2}}
          = \qty{34.7}{N}
        \end{align*}
        Answer: \final{$F = \qty{35}{N}$} (repulsive).
  \item Strengths: it explains all three gold-foil results (most straight
        through, some small deflections, a very few turned back). Geiger and
        Marsden's 1913 counts at many angles matched its predictions from
        Coulomb's law. So the evidence for a tiny, charged nucleus is strong.
        Limitations: an orbiting electron is accelerating, so classical physics
        says it should give out light and spiral into the nucleus. The model
        also cannot explain line spectra, or what makes up the nucleus.
        Judgement: the model was a major step and is strongly supported for the
        nucleus, but it is incomplete for the electrons. Bohr's model
        (\cref{ch3:chap}) improved it.
  % ---- 2.4 neutron
  \item Mass numbers: $4 + 14 = 17 + A$, so $A = 1$. Atomic numbers:
        $2 + 7 = 8 + Z$, so $Z = 1$. The particle is \final{$\chTwoNuc{1}{1}{H}$},
        a proton (also written $\chTwoNuc{1}{1}{p}$).
  \item Paraffin wax contains a lot of hydrogen, so it is full of protons.
        A proton has about the same mass as the unknown particle. So in a
        head-on collision it can take almost all of the particle's energy. The
        knocked-out protons are charged, so they ionise the gas in the detector,
        and their energy can be measured. From these energies Chadwick worked
        out the properties of the neutral radiation.
  \item (a) Choose right as positive. Momentum before $=$ momentum after:
        $m_\mathrm{n}u_\mathrm{n} = m_\mathrm{n}v_\mathrm{n} + m_\mathrm{N}v_\mathrm{N}$.
        \begin{align*}
          m_\mathrm{n}u_\mathrm{n} &= \qty{1.675e-27}{kg}\times\qty{3.00e7}{m.s^{-1}}
            = \qty{5.025e-20}{kg.m.s^{-1}}\\
          m_\mathrm{N}v_\mathrm{N} &= \qty{2.325e-26}{kg}\times\qty{4.03e6}{m.s^{-1}}
            = \qty{9.370e-20}{kg.m.s^{-1}}
        \end{align*}
        Subtract $m_\mathrm{N}v_\mathrm{N}$ from both sides, then divide by $m_\mathrm{n}$:
        \begin{align*}
          v_\mathrm{n} &= \frac{m_\mathrm{n}u_\mathrm{n} - m_\mathrm{N}v_\mathrm{N}}{m_\mathrm{n}}
          = \frac{\qty{5.025e-20}{kg.m.s^{-1}} - \qty{9.370e-20}{kg.m.s^{-1}}}{\qty{1.675e-27}{kg}}\\
          &= -\qty{2.59e7}{m.s^{-1}}
        \end{align*}
        Answer: \final{$v_\mathrm{n} = \qty{2.59e7}{m.s^{-1}}$ backwards} (the
        neutron bounces back).\par
        (b) Compare the nitrogen's kinetic energy with the neutron's starting energy:
        \begin{align*}
          K_\mathrm{before} &= \tfrac12\times\qty{1.675e-27}{kg}\times(\qty{3.00e7}{m.s^{-1}})^2 = \qty{7.54e-13}{J}\\
          K_\mathrm{N} &= \tfrac12\times\qty{2.325e-26}{kg}\times(\qty{4.03e6}{m.s^{-1}})^2 = \qty{1.89e-13}{J}\\
          \text{fraction} &= \frac{\qty{1.89e-13}{J}}{\qty{7.54e-13}{J}} = 0.25
        \end{align*}
        Answer: \final{about a quarter} (25\%), compared with almost 100\% for a proton target.
  \item Suppose the radiation were gamma rays. A photon would need about
        \qty{50}{MeV} to give a proton its measured \qty{5.7}{MeV}. But if
        beryllium captured an alpha and gave out a gamma ray, the photon could
        carry at most about \qty{14}{MeV}. So gamma rays would break conservation
        of energy. Also, no single gamma-ray energy fitted the recoils of both
        hydrogen and nitrogen. A neutral particle with a proton's mass explains
        all the results.
\end{chapteranswers}
